\section{Introduction}

Attention is usually written as
\begin{equation}
  S=QK^\mathsf{T}/\sqrt d,\qquad
  P=\softmax(S),\qquad
  O=PV.
  \label{eq:attention-intro}
\end{equation}
The fastest GPU implementation is determined not only by these equations but
also by where each tile lives and when it becomes ready. FlashAttention avoids
writing the quadratic $S$ and $P$ matrices to high-bandwidth memory
\cite{dao2022flashattention}. Later versions improved work partitioning and
adapted the pipeline to asynchronous matrix units
\cite{dao2023flashattention2,shah2024flashattention3}. FlashAttention-4 (FA4)
redesigns this schedule for Blackwell's tensor memory (TMEM), TMA transfers,
and fully asynchronous tensor-core operations
\cite{zadouri2026flashattention4}.

Blackwell also provides block-scaled FP4 matrix multiplication with a peak
throughput target near four times BF16 on GB200. That ratio applies to QK and
PV, but not to the work between them. Q, K, and V can be quantized before the
attention kernel starts. P cannot: each online-softmax iteration must create,
scale, pack, and publish P before the next matrix product can consume it.
Making QK and PV shorter therefore makes this serial path more visible.

HAO AI Lab's exploratory FP4 FA4 implementation demonstrates the central
counterexample: NVFP4 QK with FP8 PV is often faster than an all-NVFP4 route
\cite{hao2026fp4flashattention,zhang2026attnqat}. Its two-query CTA and TMEM
ownership model solve the large scheduling problem. We keep that structure
and ask a narrower question: can the online probability representation be
changed so that full FP4 PV becomes worthwhile?

Our retained format split is
\begin{equation}
  \boxed{\text{NVFP4 }QK \quad + \quad \text{MXFP4 }P/V.}
  \label{eq:retained-format}
\end{equation}
NVFP4 keeps fine block-16 scales for signed Q and K. MXFP4 gives each block of
32 P and V values a power-of-two E8M0 scale. This wider range is useful for
small probability blocks and lets scale selection fold into a log-domain code
map.

This report makes four contributions:
\begin{enumerate}
  \item It reconstructs HAO's two-query full-FP4 pipeline in ThunderKittens
  and states its TMEM ownership and N32-to-K64 publication rules explicitly.

  \item It introduces an NV/MX probability path that maps log scores directly
  toward E2M1 codes with packed FMA, mixes native EX2 with arithmetic work,
  and normalizes by the probabilities that PV actually consumed.

  \item It adds a low-cost range guard for extreme model logits. The final
  version enforces one exponent budget and uses an underflow-safe but
  algebraically equivalent denominator evaluation, fixing the observed
  20-step Wan failure without a full row scan.

  \item It evaluates latency and numerical error from matched records on
  GB200 and B300, compares against HAO and BF16, and adds ViT, BERT, ViT-MAE,
  and Wan model replays.
\end{enumerate}

The scope is forward, noncausal D128 and D64 attention. Kernel timings include
QK, online softmax, P quantization, PV, correction, and the output epilogue;
they exclude dynamic Q/K/V quantization and optional K/V reordering. The
paper first explains the inherited storage schedule, then the new probability
path, followed by matched performance, accuracy, and remaining limits.
